\section{What the measured change can support}

For this archived, English, six-item comparison against an unweighted survey
derivative, GPT-5.6 reduces mean category-mass error mainly on Q106 and Q108.
The other four corrected global contrasts are inconclusive. This does not show
better advice, better representation of individuals, or greater user welfare.

The exploratory economic grouping assigns the six items to three declared
themes, with the largest estimated reduction on distribution/adjustment
($-0.0213$). These themes are not validated economic choices; the Q106 label
mismatch further limits that margin. Appendix~\ref{app:economics} gives the
weight sensitivity and distinguishes observed losses from a hypothetical
decision bridge.

The finite five-run identity isolates the reduction from averaging these
outputs; the generations do not interact. It is a pre-interaction baseline
for future agent experiments, not evidence of deliberation or collective
behavior. Archived outputs and hashes reproduce this analysis, while the
undated \texttt{gpt-5.6-sol} identifier cannot ensure identical live answers.

English country labels may reflect questionnaire familiarity or country
stereotypes and conceal within-country language, region, class, and
institutions \citep{nekoto2020participatory,mohamed2020decolonial,sen1999development}.
No community partner co-designed or locally validated this wave. An official
population-weighted WVS comparison and locally translated prompts have not
been run; their effects on $\Delta\mathrm{TVD}$, CRG, or VDR cannot be inferred
from these data. Before using a model update for consequential advice, repeat
the item-level comparison with relevant survey weights and local-language
prompts, inspect country variation, and involve affected communities.
Appendix~\ref{app:future-roadmap} outlines further work. Until those checks
exist, survey fidelity is a diagnostic on this target, not local legitimacy.
